\section{Ablations and Discussion}
\label{sec:ablations}

This section interprets the benchmark against the three paper claims. Several experiments needed for causal attribution are absent, so we separate supported observations from missing ablations.

\subsection{C1: sequence length}

The benchmark supports the directional claim that short operational windows are effective on N-CMAPSS. At 1000 timesteps, WaveNet reaches RMSE 6.523 on FD1 (T1), MSTCN reaches RMSE 7.604 in the same FD1 protocol (T3), and MSTCN reaches RMSE 6.495 on FD2 (T11). Full-flight baselines are much weaker: the full-flight TCN row reports RMSE 16.13 (T19), while full-flight LSTM and BiLSTM rows report RMSE 22.28 (T20) and 22.20 (T21). Very short 250-step inputs also degrade performance, with MSTCN at RMSE 10.748 (T14) and WaveNet at RMSE 11.169 (T15), suggesting that the useful regime is not simply ``as short as possible.''

The limitation is that no single architecture has been trained across the full length sweep $\{100,250,500,1000,2000,5000,20000\}$ under matched conditions. The 1000-step winners are mostly attention or dilated-convolution models, whereas the 20k rows are mostly recurrent or simpler convolutional baselines. The apparent 1k-vs-full gap therefore conflates sequence length with architecture family. We report C1 as a directional finding, corroborated by short-window N-CMAPSS challenge submissions~\cite{lovberg2021variable,solismartin2021stacked}, not as a controlled 58\% single-factor ablation.

\subsection{C2: multi-scale temporal attention}

MSTCN is competitive but not uniformly dominant. In the controlled FD1 table, WaveNet ranks first (T1: RMSE 6.523), CNN-GRU second (T2: 7.006), and MSTCN third (T3: 7.604). This contradicts any simple claim that MSTCN is the FD1 winner. On FD2, however, MSTCN is best by RMSE (T11: 6.495), ahead of WaveNet (T12: 6.739), and CNN-GRU degrades sharply (T13: 19.812).

The multiseed runs (Table~\ref{tab:multiseed}) further bound the single-run estimates: MSTCN's mean RMSE over three seeds is 8.140 ($\pm0.74$, T5) and WaveNet's over two seeds is 8.238 ($\pm1.44$, T6). The means are statistically indistinguishable and both lie above their respective best-run values (T3: 7.604; T1: 6.523), so the FD1 gap between WaveNet and MSTCN in the best-run table is within the variance envelope of both models.

The interpretation supported by the current results is that MSTCN and WaveNet belong to the same top cluster on FD1, MSTCN transfers well to FD2, and full-flight recurrent baselines underperform. The architectural attribution remains open. We have not yet ablated MSTCN's multi-scale branches, removed its global fusion attention, or matched parameter counts against WaveNet and Transformer. Those experiments are required before claiming that the multi-scale or fusion-attention components themselves cause the observed performance.

\subsection{C3: asymmetric loss}

The project uses asymmetric MSE as its default loss, and the controlled rows show low PHM scores and high threshold accuracy: WaveNet on FD1 reaches Accuracy@20 98.83 and PHM 0.7256 (T1), and MSTCN on FD2 reaches Accuracy@20 99.01 and PHM 0.5755 (T11). The loss is aligned with the asymmetric PHM scoring tradition~\cite{saxena2008damage} and with recent cost-sensitive RUL loss design~\cite{mpm2024loss}.

However, C3 is not yet an empirical result of this paper. There is no symmetric-MSE baseline in the current canonical results, no $\alpha$ sweep over the late-penalty ratio, and the high-accuracy recipe in the repository uses asymmetric Huber rather than asymmetric MSE (T23: RMSE 9.132, Accuracy@20 93.84). We therefore frame asymmetric MSE as a motivated methodological choice, not as a proven improvement over symmetric MSE in our experiments.

\subsection{Open experiments}

The next experiments are specific. A single-architecture sequence-length sweep on the top cluster would isolate the C1 effect. MSTCN variants with one dilation scale and with global fusion attention removed would test the C2 mechanism. Symmetric MSE, asymmetric MSE, and asymmetric Huber under identical settings, including an $\alpha \in \{1.0,1.5,2.0,3.0\}$ sweep, would turn C3 from a design choice into an empirical comparison.
